\documentclass[12pt]{article}
\usepackage[utf8]{inputenc}
\usepackage[T1]{fontenc}
\usepackage{hyperref}
\title{Tri\'o Language Resources for K-12 Curriculum Development}
\author{Compiled by Claude Sonnet 5 (claude-sonnet-5) for Andr\'as Kornai}
\date{}
\begin{document}
\maketitle

\section{Introduction}
Tri\'o (also Tiriy\'o, Trio, Tiriyo; ISO 639-3 \texttt{tri}) is a Cariban language of Suriname and Brazil. Sources describe it as spoken by the whole community, and UNESCO classes it as Vulnerable. The search found descriptive literature but no school materials.

\section{Primary Educational Resources}
The search found no Tri\'o primer, textbook or other classroom material. This section is empty.

\section{Linguistic and Reference Materials}
\begin{itemize}
\item Eithne Carlin, \emph{A Grammar of Trio: A Cariban Language of Suriname} (Peter Lang, 2004), a comprehensive descriptive grammar: \url{https://www.peterlang.com/view/title/46332?print}
\item Sergio Meira, \emph{A Grammar of Tiriy\'o} (Rice University PhD dissertation, 1999), cited in the search results, which do not give a link to the text. Reference listing: \url{https://glottolog.org/resource/reference/id/54666}
\item Phonological inventory of Trio (SAPHON, Berkeley): \url{https://linguistics.berkeley.edu/saphon/en/inv/Trio.html}
\item Bibliography for Tiriy\'o in the Endangered Languages Project: \url{https://www.endangeredlanguages.com/lang/2137/bibliography?hl=en}
\item Overview article on the language: \url{https://en.wikipedia.org/wiki/Tiriy%C3%B3_language}
\end{itemize}

\section{Religious and Community Resources}
None were found in the search results.

\section{Digital and Audio Resources}
None were found in the search results. One source mentions a preliminary English--Tiriy\'o dictionary, but the results give no author, date or link.

\section{Curriculum Development Context}
Tri\'o is spoken by the entire community, so a curriculum could be built on spoken-language use. The only resources found are academic grammars. Primers, a standard orthography for school use, and audio would have to be created or located by further work, ideally with Tri\'o communities and mission or education bodies in Suriname and Brazil.
\end{document}